\begin{derivation}[从\eqnrefs{eq:macqed-Eplus-SI-r12}到\eqnrefs{eq:macqed-vacuum-field-correlation-dp68}]
把正频场写成

\begin{equation*}
\hat E_i^{(+)}(\mathbf r,\omega)
=\int\dd^3 s\,K_{ia}(\mathbf r,\mathbf s;\omega)
\hat f_a(\mathbf s,\omega),
\end{equation*}
其中 \(K\) 包含正文\eqnrefs{eq:macqed-Eplus-SI-r12} 中的全部数值核。真空满足
\begin{equation*}
\langle0|\hat f_a(\mathbf s,\omega)
\hat f_b^\dagger(\mathbf s',\omega')|0\rangle
=\delta_{ab}\delta^{(3)}(\mathbf s-\mathbf s')\delta(\omega-\omega').
\end{equation*}
代入两份场算符后，储库的空间 \(\delta\) 函数消去一个积分：
\begin{align*}
&\langle0|\hat{\mathbf E}^{(+)}(\mathbf r,\omega)
\hat{\mathbf E}^{(-)}(\mathbf r',\omega')|0\rangle\\
&=\mu_0^2\omega^4\frac{\hbar\varepsilon_0}{\pi}
\delta(\omega-\omega')
\int\dd^3 s\,
\mathbf G^R(\mathbf r,\mathbf s;\omega)
\cdot\operatorname{Im}\boldsymbol\varepsilon_r(\mathbf s,\omega)
\cdot(\mathbf G^R)^\dagger(\mathbf r',\mathbf s;\omega).
\end{align*}
利用 \(\mu_0\varepsilon_0=1/c^2\) 和正文\eqnrefs{eq:green-optical-identity-SI-r12}，积分等于 \((c^2/\omega^2)\operatorname{Im}\mathbf G^R\)。于是
\begin{equation*}
\langle0|\hat{\mathbf E}^{(+)}(\mathbf r,\omega)
\hat{\mathbf E}^{(-)}(\mathbf r',\omega')|0\rangle
=\frac{\hbar\mu_0}{\pi}\omega^2
\operatorname{Im}\mathbf G^R(\mathbf r,\mathbf r';\omega)
\delta(\omega-\omega'),
\end{equation*}
即正文\eqnrefs{eq:macqed-vacuum-field-correlation-dp68}。
\end{derivation}

\begin{derivation}[从\eqnrefs{eq:reservoir-thermal-occupation-dp68}到\eqnrefs{eq:reservoir-thermal-correlators-dp43}]
先看一个频率为 $\omega$ 的离散玻色模式。热平衡密度算符为
\begin{equation*}
 \rho_{\rm th}
 =\frac{\ee^{-\beta\hbar\omega\hat a^\dagger\hat a}}
 {\operatorname{Tr}\ee^{-\beta\hbar\omega\hat a^\dagger\hat a}}
 =(1-q)\sum_{n=0}^{\infty}q^n|n\rangle\langle n|,
 \qquad q=\ee^{-\beta\hbar\omega}.
\end{equation*}
归一化使用几何级数 $\sum_{n\ge0}q^n=(1-q)^{-1}$。因此
\begin{align*}
 \langle\hat a^\dagger\hat a\rangle_{\rm th}
 &=(1-q)\sum_{n=0}^{\infty}nq^n\\
 &=(1-q)q\frac{\dd}{\dd q}\frac{1}{1-q}
 =\frac{q}{1-q}
 =\frac{1}{\ee^{\beta\hbar\omega}-1}
 \equiv n_B(\omega;T_{\rm th}).
\end{align*}
相反次序不能简单把两个算符交换；由
$\hat a\hat a^\dagger=\hat a^\dagger\hat a+I$，
\begin{equation*}
 \langle\hat a\hat a^\dagger\rangle_{\rm th}
 =n_B+1.
\end{equation*}
同一个结果也可直接从数态看出：
$\hat a\hat a^\dagger|n\rangle=(n+1)|n\rangle$。而
$\langle\hat a\hat a\rangle=\langle\hat a^\dagger\hat a^\dagger\rangle=0$，因为热态在数基中对角，而这两个算符把 $|n\rangle$ 移到与其正交的 $|n\mp2\rangle$。

宏观 QED 储库只是把这个单模结果推广到连续标签
$\lambda=(i,\mathbf r,\omega)$。规范化对易关系
\begin{equation*}
 [\hat f_i(\mathbf r,\omega),
 \hat f_j^\dagger(\mathbf r',\omega')]
 =\delta_{ij}\delta^{(3)}(\mathbf r-\mathbf r')
 \delta(\omega-\omega')
\end{equation*}
说明不同连续模式彼此独立，因此离散的 Kronecker $\delta$ 被 Dirac $\delta$ 取代：
\begin{align*}
 \langle\hat f_i^\dagger(\mathbf r,\omega)
 \hat f_j(\mathbf r',\omega')\rangle
 &=n_B(\omega;T_{\rm th})\,
 \delta_{ij}\delta^{(3)}(\mathbf r-\mathbf r')\delta(\omega-\omega'),\\
 \langle\hat f_i(\mathbf r,\omega)
 \hat f_j^\dagger(\mathbf r',\omega')\rangle
 &=[n_B(\omega;T_{\rm th})+1]\,
 \delta_{ij}\delta^{(3)}(\mathbf r-\mathbf r')\delta(\omega-\omega').
\end{align*}
这就是正文\eqnrefs{eq:reservoir-thermal-correlators-dp43}。$T_{\rm th}\to0$ 时 $n_B\to0$，但第二行的单位项保留；它来自玻色对易代数本身，而不是来自热占据，因此对应真空涨落。
\end{derivation}
